% !TEX root = ../b-rg.tex
\chapter{Phonology}\label{ch:phonology}

\section{Phonemes}\label{sec:phonemes}

\subsection{Consonants}
Borlish has a consonant inventory of middling size. Notable is the presence of \phm{ʒ} without its voiceless counterpart \phm{ʃ}, due to a relatively recent sound change which backed the latter to \phm{x}. 
\begin{table}[ht]
    \centering
    \begin{tabular}{ c c c c c c }
        \multicolumn{6}{ c }{Consonants} \\
        \hline
        m & & n & & \\
        p b & & t d & & k g & \\
        & & ts & & & \\
        f v & θ ð & s z & ʒ & x & h \\
        w\textsuperscript{1} & & l & j & r\textsuperscript{2} & \\
        \hline
    \end{tabular}
\end{table}

\textsuperscript{1}In native vocabulary the phoneme \phm{w} appears only after \phm{k}. This, along with alternations like   
\glphm{voc voccar | bocq bocquar}{vɔk vɔˈkar | bɔk bɔˈkwar}{\obs summons {to summon} | fee {to charge}}
have led some to analyse \phm{kw} sequences as a single phoneme \phm{kʷ} which merges with \phm{k} in coda position.

\textsuperscript{2} The rhotic is uvular, but for ease we will denote it by \phm{r} in phonemic transcription.

\subsection{Vowels}
Borlish has eight monophthong vowel phonemes, typical of neighbouring Germanic and sister Gallo-Romance languages.

\begin{table}[ht]
    \centering
    \begin{tabular}{ c c c }
        \multicolumn{3}{ c }{Vowels} \\
        \hline
        i & ɪ & u \\
        e & & o \\
        ɛ & a & ɔ \\
        \hline
    \end{tabular}
\end{table}

Borlish also contains the diphthong \phm{au}, as well as vowel+\phm{j} sequences that may be analysed as diphthongs. 

The segments \phm{m n r} can appear as syllabic consonants.
    \glphm{batesm cohernt outrscac}{baˈtɛz.m̩ koˈhɛr.n̩t ˌut.r̩ˈxak}{baptism consequent brazen}
The syllable preceding a syllabic consonant acts as a closed syllable, and the segment before a syllabic consonant acts as part of the coda of the previous syllable. Hence, it is usually considered most elegant to syllabify so that syllabic consonants have no onset.

\section{Stress}\label{sec:stress}
The final syllable is usually stressed.
    \glphm{scadom tatover}{xaˈdɔm ˌta.toˈvɛr}{tomato potato}
In words of four or more syllables, secondary stress is usually assigned in iambs, although anapaestic patterns do occur when the interior syllables are sufficiently light.
    \glphm{deborsellar haubrjoranç}{deˌbɔr.sɛˈlar ˌhob.r̩.ʒoˈrants}{pickpocket advertisement}
There are some exceptions to final stress. Syllabic consonants reject stress onto the previous syllable. Moreover, the suffixes \bl{-(e)nt} (forming present participles of \bl{-r} conjugation verbs), \bl{-(e)nç} (forming nouns from the same verbs) and \bl{-essem, -errem} (forming comparatives) also enforce penultimate stress.
    \glphm{casn pardenç pessem veðerrem}{ˈkazn̩ ˈpar.dɛnts ˈpɛ.sɛm veˈðɛ.rɛm}{oak loss worse elder}
In many dialects these suffixes are parsimoniously analysed as instead containing syllabic consonants, pronounced as \phm{-n̩t -n̩ts -ɛs.m̩ -ɛr.m̩}.

\section{Phonotactics}\label{sec:phonotactics}
There are several restrictions on where in a syllable or a word certain segments can appear. The consonants \phm{v s ʒ h} do not occur in coda position, and \phm{ð} does not occur word-finally. When morphology would place \phm{ð s ʒ} in a forbidden position they are replaced with \phm{θ z x} respectively.
    \glphm{nað-ar nað | fass-ar fas | naj-ar nascn}{naˈðar naθ | faˈsar faz | naˈʒar ˈnax.n̩ }{{swim-\Inf} swim | {wrap-\Inf} wrap | {sail-\Inf} {sail.\Third\Pl}}
    When morphology would place \phm{v} in a forbidden position, either \phm{f} replaces it or the preceding vowel is altered.
    \glphm{cav-ir caf | bav-ar bau}{kaˈvɪr kaf | baˈvar bo}{{get-\Inf} get | {bark-\Inf} bark}
Syllabic consonants do not occur in stressed or initial syllables, and the `lax' vowels \phm{ɪ ɛ ɔ} do not occur in stressed open syllables.

\subsection{Onsets}
Any vowel can begin an initial syllable, but under most analyses non-initial syllables must begin with a consonant.
    \glphm{al enç aiç id eir onc aust ou}{al ɛnts ets ɪd ir ɔnk ost u}{wing start ease ides go hook east egg}
Any single consonant may constitute an onset. Two-consonant onsets are of three kinds: for the first, the second segment is a liquid (one of \phm{l j r}; the case of \phm{w} is dealt with above), preceded either by a stop or a fricative. The liquid \phm{l} does not occur after the segments \phm{t d θ ð ʒ}, and the liquid \phm{r} does not occur after the segments \phm{ð s z ʒ}.
    \glphm{plait droug gien thron}{plet druj ʒjɛn θrɔn}{pleased contribution cheek throne}
The second class of two-consonant onsets are a heterorganic stop followed by one of \phm{s θ t d}; these usually derive from Greek.
    \glphm{xyron chthonic ptyssoscia bdella}{ksiˈrɔn kθɔˈnɪk ˌpti.soˈxja bdɛˈla}{scalpel subterraneous glaucoma virus}
The remaining two-consonant onsets comprise \phm{s} followed by a nasal, a stop, the two liquids \phm{l j} or the two fricatives \phm{f v}.
    \glphm{smout skore sdegn slanc siesc sforç svam}{smut skoˈre sdɪjn slank sjɛx sfɔrts svam}{riot urgent scorn thin siege strain sponge}
All the possible three-consonant onsets are formed by the combination of \phm{s}+consonant and consonant+liquid onsets.
    \glphm{splorar spien sdruçol sfrey}{sploˈrar spjɛn sdriˈtsɔl sfri}{explore hope lubrication horror}
No more complex onsets are permitted.
\subsection{Codas}
Empty codas are permitted, including word-finally. However, the `lax' vowels \phm{ɪ ɛ ɔ} do not occur in final open syllables (as these would necessarily be stressed).
    \glphm{idone nijomm-au}{ˌi.doˈne ˌni.ʒɔˈmo}{suitable {bind-\Pst}}
Any single consonant except \phm{v s ʒ h} may be a coda. There are several types of two-consonant coda. A liquid may be followed by any non-liquid that can appear in codas at all; a nasal by any such non-liquid and non-nasal.
    \glphm{sculd tragç corf lucern | hamt coyenç ans}{xɪld trɛjts kɔrf liˈtsɛrn | hamt koˈjɛnts anz}{obligation clue raven lantern | live prevention handle}
The fricative \phm{s} may precede a voiceless stop; \phm{f} may precede \phm{t} specifically.
    \glphm{cosp fest masq | haft}{kɔsp fɛst mask | haft}{{bee sting} banquet disguise | busy}
Conversely, \phm{s} may also follow any voiceless stop; after \phm{t} this is identical to the affricate phoneme \phm{ts}.
    \glphm{laps ouç jex}{laps uts ʒɛks}{mistake silence bullseye}
Finally among two-consonant codas, the cluster \phm{kt} is also permitted.
    \glphm{abject}{abˈʒɛkt}{despicable}
The few three-consonant codas that appear are a nasal or \phm{j} with a licit two-consonant cluster.
    \glphm{boist deunx defonct}{bɔjst daunks deˈfɔnkt}{box {near miss} deceased}
No more complex codas are permitted.

\section{Allophony}\label{sec:allophony}
Throughout this section, the phonetic transcription approximately reflects the mesolectal speech in the capital city of Damvath.

\subsection{Assimilation}
The phoneme \phm{n} is pronounced \pht{ŋ} before velars.
    \glpht{blanc viking mensc}{blank viˈkɪng mɛnx}{blaŋk vɪˈkɪŋg mɛŋx}{white raid honour}
Regressive voicing assimilation affects \phm{s} before any voiced obstruent, nasal or \phm{l}.
    \glpht{sboc svanið pasnagl desjeun bruslaç}{sbɔk svaˈnɪθ paˈsnɛjl dɛsˈʒaun briˈslats}{zbɔk zvɐˈnɪθ pɐˈznɛjl dɪzˈʝaun bʀɪˈzlats}{smirk unconscious parsnip breakfast sunburn}

\subsection{Palatals}
The phoneme \phm{x} surfaces as \pht{ç} when adjacent to a front vowel.
    \glpht{scið esc}{xɪθ ɛx}{çɪθ ɛç}{possible bait}
The phoneme \phm{ʒ}, along with \phm{j} either word-initially or between an unstressed and a stressed vowel (in that order), is pronounced \pht{ʝ}.
    \glpht{jast yon çuyal}{ʒast jɔn tsaˈjal}{ʝast ʝɔn tsɐˈʝal}{zinc where cicada}
    
\subsection{Non-rhoticity}
Syllabic \phm{r̩} surfaces as \pht{ɐ}. Furthermore, for some speakers syllabic \phm{m n} surface as \pht{ɐm ɐn} respectively.
    \glpht{issambr ðesmbal visn}{ɪˈsamb.r̩ ˌðɛz.m̩ˈbal ˈvɪz.n̩}{ɪˈsam.bɐ ˌðɛ.zɐmˈbal ˈvɪ.zɐn}{massacre dodgeball knot}
Coda \phm{r} surfaces as length on a preceding monophthong, or as \pht{ɐ} after a diphthong or \phm{j}. The vowel's quality may also be altered, with \phm{ar ir ur} becoming \pht{ɐː ɪː ʊː} respectively.
    \glpht{vert hagr tourm}{vɛrt ˈhɛj.r̩ turm}{vɛːt ˈhɛ.jɐ tʊːm}{green linchpin squadron}
When a word beginning with a vowel immediately follows, the \pht{ʀ} is resurrected as the onset of the next syllable.
    \glpht{l'hour uncos}{lur ɪnˈkɔz}{lʊː‿ ʀɪŋˈkɔz}{{\Def=time} something}

\subsection{Reduction}
In syllables without primary or secondary stress, the vowel \phm{a} surfaces as \pht{ɐ}, the vowels \phm{ɛ e i} as \pht{ɪ} and the vowels \phm{ɔ o u} as \pht{ʊ}.
    \glpht{tovaresc taisson torvegl}{ˌto.vaˈrɛx teˈsɔn tɔrˈvɪjl}{ˌto.vɐˈʀɛç tɪˈsɔn tʊːˈvɪjl}{homosexual badger whirlwind}
Before an obstruent or after \phm{n}, a coda \phm{ts} is pronounced \pht{s}.
    \glpht{amaçgat fenç}{ˌa.matsˈgat fɛnts}{ˌa.mɐsˈgat fɛns}{nursery split}
The sequence \phm{kts} is realised as \pht{ks}.
    \glpht{acceir facçon}{akˈtsir fakˈtsɔn}{ɐkˈsɪː fɐkˈsɔn}{approach party}

% --- Planned sections (from the Lingua questionnaire, 3.1.2.3 and 3.3.4) ---
% \section{Loanword phonology}
%   Segments and clusters found only in loans (w outside kw, the Greek
%   onsets xyron/chthonic/ptyssoscia, æ œ ï); how loans are nativised.
% \section{Intonation}
%   Statements vs questions (with and without final ou); emphatic and
%   contrastive stress; phrasing, and the linking [ʀ] across word
%   boundaries.
